% Located, Not Forced -- arXiv/Zenodo submission source.
% v2.13 (2026-07-23): publication-hardening release reconcile. Added the explicit class-C
%   IN/OUT boundary table, reconciled the function-space residual wording to distinguish
%   canon-grade scope from the exploration-grade true-Y14 narrowing, updated creator metadata,
%   and made the source Tectonic-compatible without changing any scientific verdict.
% v2.11 (2026-07-07): Section 9 (sec:open) sharpened -- both arenas (cohomology + framed bordism) mapped and
%   walled by the 2026-07-07 RS-index and framed-bordism swings; the open bridge collapses onto one
%   integer-by-construction object. Verdict UNCHANGED (located, not forced); no grade/claim/title moved.
% RECONCILE (2026-07-07): .tex<->.md content verified IN SYNC. The version-number lag (.tex header said v2.9)
%   was mostly a stale header; the body already carried the v2.10 content -- the status-table rows (R1 actual-RS
%   residual, R2 mod-3 arena empty, the table-free parity backstop for Theorem 1), the framed-bordism Section 8,
%   and the Data Availability section. Section 9 is now synced to v2.11 with macros aligned (\pis, \RP, \Lamp,
%   \eR). LaTeX integrity checked: environments balanced (22/22), every \ref/\cite/\label resolves, amssymb loaded.
% REMAINING before submission (Joe-side only, mechanical): compile on Overleaf to confirm rendering; then the
%   arXiv endorsement check and the explicit go. No content reconcile outstanding.
% v2.9: folded four staged 2026-07-02 findings (external-by-structure synthesis, core-theorems symbolic proof, external index=flux, RS function-space framework SPEC) into caveats and grade language.
% Draft for external review. This paper does NOT claim three generations.
%
% Suggested arXiv classification:
%   Primary:   hep-th
%   Secondary: math-ph, math.AT
% Keywords: generation number, family puzzle, Rarita-Schwinger, primary decomposition,
%           Chinese Remainder Theorem, index theorem, framed bordism, Adams e-invariant,
%           anomaly inflow, chirality.
% Submission comment (non-promotional): "Draft for external review. Reconstruction-grade
%   GU-specific steps are explicitly marked; the theorem-grade results are GU-independent."

\documentclass[11pt]{article}

\usepackage[margin=1in]{geometry}
\usepackage{amsmath,amssymb,amsthm}
\usepackage{mathtools}
\usepackage{enumitem}
\usepackage[hidelinks]{hyperref}
\usepackage{xurl}
\usepackage{booktabs}
\usepackage{longtable}
\usepackage{array}
\setlength{\emergencystretch}{3em}

\newtheorem{theorem}{Theorem}
\newtheorem{corollary}{Corollary}
\newtheorem{lemma}{Lemma}
\theoremstyle{definition}
\newtheorem{definition}{Definition}
\theoremstyle{remark}
\newtheorem*{remark}{Remark}

\newcommand{\Z}{\mathbb{Z}}
\newcommand{\R}{\mathbb{R}}
\newcommand{\RP}{\mathbb{RP}}
\newcommand{\pis}{\pi_3^{s}}
\newcommand{\Cl}{\mathrm{Cl}}
\newcommand{\Spin}{\mathrm{Spin}}
\newcommand{\SU}{\mathrm{SU}}
\newcommand{\SO}{\mathrm{SO}}
\newcommand{\Hom}{\mathrm{Hom}}
\newcommand{\id}{\mathrm{id}}
\newcommand{\Ahat}{\hat{A}}
\newcommand{\ch}{\mathrm{ch}}
\newcommand{\Lamp}{\Lambda^2_{+}}
\newcommand{\Jq}{J_{\mathrm{quat}}}
\newcommand{\eR}{e_{\R}}
\newcommand{\Imm}{\operatorname{Im}}

\title{Located, Not Forced:\\
Two-Primary Obstructions Cannot Force the Fermion Generation Count\\
in a Clifford Rarita--Schwinger Sector}

\author{Joseph Hernandez\\Independent Researcher\thanks{Correspondence: \texttt{joe@disruptionjoe.com}. This work was produced by an
AI-directed research process; all load-bearing computations are reproducible from the public repository linked
in the Data Availability section.}}

\date{July 2026}
\hypersetup{
  pdftitle={Located, Not Forced: Two-Primary Obstructions Cannot Force the Fermion Generation Count in a Clifford Rarita-Schwinger Sector},
  pdfauthor={Joseph Hernandez},
  pdfsubject={Generation count, chirality, and two-primary obstructions in a Clifford Rarita-Schwinger sector},
  pdfkeywords={generation count, chirality, Clifford algebra, Rarita-Schwinger, stable homotopy, bordism, anomaly, Krein}
}

\begin{document}

\maketitle

\begin{center}
\fbox{\parbox{0.92\textwidth}{\textbf{Draft for external review. This paper does \emph{not} claim three
generations.} The theorem-grade results are independent of Geometric Unity; every Geometric-Unity-specific step
is explicitly marked as motivated or reconstruction-grade. Five caveats are load-bearing and hold throughout:
\emph{(a)} the torsion-count reading is a \emph{premise, not a result}---under a literal integer-index reading the
same obstructions would instead \emph{forbid} an odd count outright; \emph{(b)} the projected order-$3$ component cannot
literally \emph{be} an integer count, since $\Hom(\Z/3,\Z)=0$ (no homomorphism from a torsion group to a
torsion-free one is nonzero); \emph{(c)} the class-C result is complete only for its encoded equivariant
finite carrier data and listed controls; non-equivariant operators, external backgrounds, and the true
$Y^{14}$ source-action pushforward remain open, so externality is not proved; \emph{(d)} Theorem~\ref{thm:2}
is a finite Krein intersection-nullity theorem, not a physical handedness count, Hilbert-space positivity
claim, or Fredholm index. At canon grade, faithful stand-ins discharge the conditional APS/end/gap/family
terms; transfer to the true Rarita--Schwinger $Y^{14}$ bundle remains open; \emph{(e)} every verification reported here is \emph{internal}---computations reproduced from
scratch and adversarially reviewed within the same AI-directed process that produced them; no result has yet been
independently replicated or peer-reviewed. See Sections~\ref{sec:intro} and~\ref{sec:open}.}}
\end{center}

\begin{abstract}
We study the chirality and generation structure of an explicit real Clifford Rarita--Schwinger sector: the
gamma-traceless rank-$3/2$ field of a real Clifford module $\Cl(p,q)$ with $p+q=14$ on a $128$-dimensional
spinor, with one generation identified with the $\mathbf{16}$ of $\Spin(10)$. The sector is \emph{motivated by},
and reconstructed from, the matter proposal of Geometric Unity (GU); every GU-specific step is marked, and none
of the theorem-grade results depends on GU being correct.

We prove a codomain-separated \emph{two-primary constraint theorem}: in an encoded delimited class C of
covariant, sector-interior structures, the finite torsion constraints are $2$-primary and the integer-valued
rows carry only parity/divisibility content; none supplies a mod-$3$ congruence or selects a particular odd
integer. Listed first-step departures outside C are adversarial controls, not an exhaustive classification
outside C. We also prove a finite Krein intersection theorem: a maximal positive subspace of the
cross-chirality $(96,96)$ carrier intersects each chiral Lagrangian trivially. This does not count the physical
left-minus-right handedness of its vectors and is not a Fredholm index.

The standard decomposition $\pis=\Z/24=\Z/8\oplus\Z/3$ separates torsion data already mapped into the two
summands. It does not move the paper's $\Z$-valued constraints into $\Z/8$ or an integer generation count into
$\Z/3$; those maps are absent. We
exhibit a framed class with a nonzero $3$-primary component in that summand---the framed-bordism Adams
$e$-invariant $\eR=1/12$ carried by the self-dual $\Lamp$ tangential framing on the $\RP^3$ spine of the metric
fiber. The full class $2\in\Z/24$ and $\eR=1/12\in\mathbb{Q}/\Z$ have order $12$; only the projected
$\Z/3$ component has order $3$. We show by computation that it is tangential but vectorlike and homotopy-fixed:
it \emph{locates} the $3$-primary slot without filling it
(``locates'' is an interpretive gloss under the torsion-count premise of caveat~\emph{(a)}, not a further
computed quantity; by $\Hom(\Z/3,\Z)=0$ the located class cannot itself be a count).

The resulting no-go is deliberately scoped: none of the encoded class-C constructions or listed controls
supplies a typed integer generation count. We do not infer that the count is necessarily external, because
non-equivariant operators and the true Rarita--Schwinger $Y^{14}$ bundle/source-action pushforward remain open.
External flux backgrounds provide a positive control showing how an integer index can arise, not a proof that
they are the unique route.

We do \emph{not} claim three generations. The ``located, not forced'' reading is contingent on treating the
obstruction data torsion-theoretically; under a literal integer-valued chiral-index reading the same
obstructions can instead \emph{forbid} an odd count. We claim a scoped no-go that names the missing typed
bridge; the identification
$3\text{-primary torsion component}\to\text{integer }3$ is posed as the single named open conjecture (a
candidate category error, made algebraically explicit by $\Hom(\Z/3,\Z)=0$). Primary decomposition and
positive uses of its $3$-primary content are prior art. The search-bounded candidate delta is only the
carrier-specific inverse assembly. External review of both the core results and the open bridge is requested.
\end{abstract}

\section{Introduction}\label{sec:intro}

The Standard Model contains three generations of fermions and does not explain the number. A recurring hope in
unification is that the number is fixed by topology. The obstacle, made precise here, is that the natural index-
and anomaly-theoretic invariants one computes \emph{in this sector} are predominantly \textbf{$2$-primary}:
mod-$2$ Witten anomalies, mod-$16$ topological-superconductor classifications, even-valued Dirac indices. A
$2$-primary torsion invariant cannot select a $3$-primary class through the group structure alone; an
integer-valued zero or evenness constraint can instead forbid an odd count. (Where the literature \emph{does} obtain an odd
constraint---Garc\'ia-Etxebarria--Montero's $\Z_9 \to N_{\mathrm{gen}}\in 3\Z$ is the sharpest example---the
constraining anomaly is itself genuinely $3$-primary, consistent with the thesis that an odd count requires
$3$-primary input.)

This paper turns that obstacle into a positive structural statement, in one explicit sector, and then tests the
strongest possible reading of it by direct computation. Two distinctions are load-bearing and we keep them
explicit throughout. First, \emph{motivated} versus \emph{independent}: the sector is motivated by GU, but the
theorem-grade results (the codomain-separated enumeration, the finite Krein intersection theorem, the CRT decomposition,
the $\Hom(\Z/3,\Z)=0$ sharpening) hold for the Clifford--Rarita--Schwinger sector as a piece of mathematics,
whether or not GU is correct. Second, \emph{located} versus \emph{forced}: a nonzero $3$-primary torsion
projection is located, but no typed integer generation count or class-to-count map is constructed. Under a
literal integer-valued reading some constraints instead forbid an odd count. We keep those codomains separate.

Our contributions:
\begin{enumerate}[leftmargin=2em]
\item \textbf{A codomain-separated two-primary constraint theorem (Section~\ref{sec:thm1}).} No item in the
encoded class-C list supplies an odd-prime congruence or selects a particular odd integer. Finite torsion rows
are $2$-primary; integer rows remain $\Z$-valued parity/divisibility constraints. Listed outside-C controls
are not an exhaustive classification beyond C.
\item \textbf{A CRT two-arena reading (Section~\ref{sec:crt}).} The decomposition
$\pis = \Z/24 = \Z/8 \oplus \Z/3$ is standard primary decomposition, already used in the prior family-puzzle
literature; we claim no novelty for it. If an explicit map places a finite sector constraint in $\Z/8$, that
image cannot determine the separate $\Z/3$ component. No such statement is made for the $\Z$-valued rows or
for an integer generation count.
\item \textbf{A finite Krein intersection theorem and bounded operator audit
(Section~\ref{sec:thm2}).} Every maximal $K$-positive subspace intersects each chiral Lagrangian trivially, and
linear $K$-isometries preserve that nullity. This is not a physical handedness count. The same intersection
statement holds for a delimited class of antilinear $K$-null re-gradings. No encoded class-C construction
supplies the missing integer generation operator; unrestricted and true-$Y^{14}$ constructions remain open.
\item \textbf{A located $3$-primary carrier (Section~\ref{sec:carrier}).} The self-dual $\Lamp$ tangential
framing on the $\RP^3$ spine carries $\eR=1/12$. The full class has order $12$, while its projected $\Z/3$
component is nonzero and has order $3$. It is shown by computation to be tangential but vectorlike and
homotopy-fixed: it locates, it does not fill (``locates'' is an interpretive gloss under the torsion-count
reading, not an additional computed quantity).
\item \textbf{The deciding computation and the gate (Section~\ref{sec:decider}).} The integer count on GU's
actual $14$-manifold is gated on the unbuilt stabilized twisted Rarita--Schwinger source action; the only
unconditionally computable integer in this program is one. We do not force three.
\item \textbf{The open conjecture (Section~\ref{sec:open}).}
$3\text{-primary torsion component}\to\text{integer }3$ is,
in the present literature, a theorem of nothing and a candidate category error, made precise by
$\Hom(\Z/3,\Z)=0$. We state it as the single open bridge and request review.
\end{enumerate}

The robust, theory-independent content is the codomain audit, finite intersection theorem, standard CRT
decomposition, and framed-class calculation. The GU
identification and the generation-count reading are clearly marked reconstruction-grade or open.

\section{The setting}\label{sec:setting}

Let $S$ be the real spinor module of $\Cl(p,q)$, $p+q=14$, $\dim_\R S = 128$. The Rarita--Schwinger field is the
gamma-traceless part of $V\otimes S$; the constraint cuts $V\otimes S$ (dimension $14\cdot128=1792$) to
$\ker(\Gamma)$, of dimension $(14-1)\cdot128 = 1664 = 2^7\cdot 13$. We split $14 = 4+10$: a four-dimensional base
$X^4$ and a ten-dimensional internal space, with one generation the $\mathbf{16}$ of $\Spin(10)$. The base frame
group is $\Spin(4)=\SU(2)_+\times\SU(2)_-$; the self-dual two-forms $\Lamp$ generate $\mathfrak{su}(2)_+$. Under
$\mathfrak{su}(2)_+$,
\[
\ker(\Gamma) = 640~(\text{singlets}) + 832~(\text{doublets}) + 192~(\text{triplets}),
\]
verified numerically. The $192$-dimensional $j=1$ sector is the native multiplicity-three structure: a genuine
triplet carrying the $\mathbf{16}$. The invariant bilinear is the $\mathfrak{so}(p,q)$-invariant Krein form
$K = \eta_V\otimes\beta_S$; on the triplet it is purely cross-chirality, signature exactly $(+96,-96)$, so the
matter content is \emph{vectorlike} with net chirality $0$. This is the central tension: the sector is
multiplicity-three but vectorlike, and no internal operation chiralizes it.

\section{Multiplicity versus chirality}\label{sec:multiplicity}

The integer $3$ plays two arithmetically different roles in this sector, and conflating them is the error the
``located, not forced'' reading most guards against. A \emph{multiplicity} is the dimension of a flat family
space---the number of copies of a fixed representation the constrained module carries; it is a representation
dimension, comes with its mirror (copies and anti-copies), and is vectorlike. A \emph{net chiral count} is an
index---the signed difference (left minus right) of Dirac zero modes; it measures chirality, the property a
vectorlike spectrum lacks. A representation dimension and an operator index are not one number computed two ways.
This sector supplies the first natively and does not supply the second.

\paragraph{The multiplicity is natively three (verified).} Under the self-dual $\mathfrak{su}(2)_+$ of the
$4$-base, the gamma-traceless module decomposes as in Section~\ref{sec:setting}. Written by \emph{multiplicity}
rather than dimension,
\[
\ker(\Gamma) = 640~(j{=}0) + 416~(j{=}\tfrac12) + 64~(j{=}1),
\qquad 640\cdot 1 + 416\cdot 2 + 64\cdot 3 = 1664,
\]
i.e.\ $640$ singlets, $416$ doublets and $64$ triplets, the $64$ triplets spanning the $192$-dimensional $j=1$
sector of Section~\ref{sec:setting}. On all $192$ states the $\Spin(10)$ Casimir matches the reference
$\mathbf{16}/\overline{\mathbf{16}}$ value to machine precision
(\path{tests/generation-sector/h1_selfdual_family_kill.py}). A native multiplicity-three object exists,
carried by the self-dual geometry of any oriented $4$-base; the multiplicity dimension $3$ is not imported.

\begin{lemma}[spinor $2$-smoothness]\label{lem:2smooth}
If a family symmetry $\SO(m)$ acts so that the multiplicity space is a sum of spinor representations of $\SO(m)$,
its dimension is a power of two, hence not divisible by $3$. An odd multiplicity can therefore arise only from a
non-spinor family representation---a vector, an adjoint, or a self-dual tensor.
\end{lemma}
\begin{proof}
The Dirac spinor of $\SO(m)$ has dimension $2^{\lfloor m/2\rfloor}$ and each half-spinor of $\SO(2k)$ has
dimension $2^{k-1}$; all are powers of two, never divisible by the odd prime $3$. The vector ($m$), the adjoint
($m(m-1)/2$) and self-dual tensors are not so constrained; $\Lamp$ of a $4$-manifold has dimension $3$.
\end{proof}

\paragraph{Three is the canonical odd multiplicity (verified for the commutant of the full $\Spin(10)$).}
Lemma~\ref{lem:2smooth} turns an open-ended search into a dichotomy: the spinor channels are provably $3$-free,
so any odd multiplicity must enter through a non-spinor channel. Inside $\Spin(14)\supset\Spin(10)\times\Spin(4)$
the commutant of $\Spin(10)$ is $\mathfrak{so}(4)=\mathfrak{su}(2)_+\oplus\mathfrak{su}(2)_-$, which up to
conjugacy has exactly three $\mathfrak{su}(2)$ subalgebras (self-dual, anti-self-dual, diagonal), and no
$\mathfrak{su}(3)$ embeds (by dimension, $\dim\mathfrak{su}(3)=8>6=\dim\mathfrak{so}(4)$). Enumerating the
$16$-dimensional multiplicity space under each
(\path{tests/generation-sector/leg3_family_embedding_enumeration.py}), only the self-dual and anti-self-dual
$\mathfrak{su}(2)$ (orientation images of one another) give an odd multiplicity, always the triplet of dimension
$3$, never a quintet or higher; the diagonal choice gives an all-even content. So for family symmetries commuting
with the full $\Spin(10)$ the self-dual route is canonical and the multiplicity it fixes is three. \emph{This is a
statement about a representation dimension, not a net chiral count.} The case of a family symmetry commuting only
with the Standard-Model subgroup $G_{\mathrm{SM}}\subset\Spin(10)$---whose larger commutant could in principle
host a horizontal $\mathfrak{su}(3)$---is not covered here and is a bounded open computation; Lemma~\ref{lem:2smooth}
already constrains it.

\paragraph{But the triplet does not by itself supply a chiral count.} The invariant Krein form
$K=\eta_V\otimes\beta_S$ on the triplet is purely cross-chirality with signature $(+96,-96)$.
Theorem~\ref{thm:2} proves only a finite Krein-side intersection nullity: a maximal $K$-positive subspace
contains no nonzero vector lying wholly in either chiral Lagrangian. It does not count the handedness of the
$96$ physical vectors, because a generic graph vector has components in both chiral spaces. Turning the native
multiplicity three into a net physical three therefore requires a separately typed grading-invariant subspace,
graded trace, or Fredholm/index construction, none of which is supplied by this finite carrier argument.
\textbf{The multiplicity three is native; a physical chiral count is not derived here.}

\section{Theorem 1: codomain-separated two-primary constraints}\label{sec:thm1}

A finite torsion invariant is \textbf{$2$-primary} when it is annihilated by a power of two. An
integer-valued constraint is called \textbf{$2$-adic/parity-valued} here only when its content is divisibility
by a power of two. These are different codomains: a $\Z$-valued equality or index is not an element of
$\Z/8$ unless an explicit reduction or torsion-valued map is supplied.

\begin{theorem}\label{thm:1}
Among the constraints on a candidate generation count enumerated for this sector (items~(1)--(7) in the proof
below), none supplies an odd-prime congruence or selects a particular odd integer. The finite torsion
constraints are $2$-primary; the integer-valued equalities and indices are parity/divisibility constraints and
remain $\Z$-valued. In particular, no enumerated item produces a mod-$3$ condition or a value-three selection.
The enumeration is complete for the encoded delimited class C of
covariant, sector-interior structures (linear or antilinear on the carrier, equivariant under the sector's split
symmetry $\mathrm{so}(4)\oplus\mathrm{so}(10)$, built from the sector's own data): a computed-grade classification
finds no class-C generator with odd-prime congruence content. Listed first-step departures outside C
(gauge-twist, boundary/$\eta$, composition, and dressed-pairing controls) were tested but are not an exhaustive
classification outside C; odd primes appear after importing OUT-side data ($54\to3$, $120\to7$,
$126\to5\cdot7$). Completeness over the full unrestricted operator theory---external
backgrounds, gauge twists by reps the sector does not contain, and the function-space / Fredholm setting---remains
open; the claim is not an impossibility statement over all conceivable obstructions.
\end{theorem}

\begin{center}
\small
\begin{tabular}{p{0.18\textwidth} p{0.35\textwidth} p{0.39\textwidth}}
\toprule
\textbf{Boundary} & \textbf{IN class C} & \textbf{OUT of class C} \\
\midrule
Carrier & Fixed $192$-dimensional $j=1$ Clifford--RS generation carrier & Added fields, representations, or spurions not present in the carrier \\
Maps & Linear or antilinear carrier maps equivariant under $\mathfrak{so}(4)\oplus\mathfrak{so}(10)$ & Non-equivariant maps or an added global/discrete symmetry such as $\Z/9$ \\
Data & Invariant forms, traces, graded traces, Kramers signs, sector index data, and characteristic classes from the sector's own bundles & External backgrounds, non-sector gauge twists, or background vacuum expectation values \\
Analytic setting & Finite carrier and sector-interior algebraic/index data & Full section-space, Fredholm, APS/family-index, and true-$Y^{14}$ source-action data \\
\bottomrule
\end{tabular}
\end{center}

\noindent
``OUT'' means outside the theorem's quantified class, not impossible or physically inadmissible. The extension
engine tests several first-step departures from C as adversarial controls; odd-primary factors appear only
after importing data on the OUT side of this boundary.

\begin{proof}[Proof (enumeration and codomain audit)]
We enumerate the constraints and keep their codomains explicit.
\begin{enumerate}[label=(\arabic*),leftmargin=2.2em]
\item \emph{Kramers / quaternionic wall.} The quaternionic structure $J^2=-1$ is a $\Z/2$ statement (Kramers
degeneracy pairs states two at a time).
\item \emph{Real / pseudoreal non-chirality.} The relevant Witten index is the mod-$2$ index, $\Z/2$-valued; a
real or pseudoreal representation cannot carry a net chirality detectable beyond mod $2$.
\item \emph{Cross-chirality Krein signature.} The $(+96,-96)$ split yields the finite intersection-index
nullity of Theorem~\ref{thm:2}. This is a $\Z$-valued equality, not a torsion class and not a physical
left-minus-right state count.
\item \emph{Adjoint Dirac index.} The adjoint index $2\,T(\mathrm{adj})\,k = 4k$ is divisible by $4$; over the
full multiplicity bundle it is $12k$ or $24k$. These are $\Z$-valued divisibility statements, not elements
of $\Z/8$.
\item \emph{Rokhlin}~\cite{rokhlin1952}. For a spin $4$-manifold, $\mathrm{sign}(X)\equiv 0 \pmod{16}$, a
mod-$2^4$ statement.
\item \emph{Spinor $2$-smoothness.} A spinor of $\SO(m)$ has dimension $2^{\lfloor m/2\rfloor}$, a power of two;
only a non-spinor (vector, adjoint, or self-dual) family representation can be odd-dimensional.
\item \emph{Ghost parity.} The tested ghost-parity resolution of the Krein hyperbolic pairs yields a $50/50$
balance diagnostic. This is again an integer/rank statement, not a torsion class.
\end{enumerate}
No item supplies a mod-$3$ congruence or selects the integer three. The finite torsion rows are genuinely
$2$-primary. The integer rows can forbid an odd value under a literal integer-index reading; they cannot be
moved into the $\Z/8$ summand without an additional map.
\end{proof}

\begin{remark}
The integer $3$ does appear in the sector---as a multiplicand ($96 = 2^5\cdot 3$, $12k$, $24k$,
$|\mathrm{Weyl}(D_7)|$). Theorem~\ref{thm:1} is not ``no factor of $3$ occurs.'' It is the sharper statement that
no enumerated constraint supplies an odd-prime \emph{congruence} or selects a specific odd integer. The
$\Z$-valued rows and the finite torsion rows must not be conflated.
\end{remark}

\begin{remark}[Structural backstop, 2026-07-03, exploration-grade]
Beyond the enumeration, the parity of any covariant sector-interior net chiral count is fixed without a census. The full chirality is central in the even Clifford algebra, so by Schur every $\mathfrak{so}(4)\oplus\mathfrak{so}(10)$-irreducible subrep of the carrier sits wholly in one chirality, and the achievable net counts form the lattice $\{\sum_R d_R(k_R-l_R)\}$; an odd count is possible iff some appearing irrep has odd dimension. On the actual carrier every appearing irrep is even-dimensional, so every covariant interior count is even---a table-free parity theorem that backstops the enumeration's 2-primary verdict, addressing the ``is the enumeration complete?'' concern structurally rather than by list (exploration-grade).
\end{remark}

\section{The CRT two-arena structure}\label{sec:crt}

We have $\pis = \Z/24$, with primary decomposition $\Z/24 = \Z/8 \oplus \Z/3$. Adams' image-of-$J$ theorem~\cite{adams1966jx4}
gives, in stem $4s-1$, that $\Imm J$ is cyclic of order $\mathrm{denom}(B_{2s}/4s)$; for $s=1$ this is $24$, so
$\Imm J_3 = \Z/24$ and the Adams $e$-invariant detects it.

\paragraph{The central structural fact and its limit.} By the Chinese Remainder Theorem the splitting
$\Z/24 = \Z/8\oplus\Z/3$ is into disjoint summands: there is no nonzero homomorphism
$\Z/8\to\Z/3$ or $\Z/3\to\Z/8$ (Lean-backed in
\texttt{canon/two-arena-rep-theory-core-RESULTS.md}). This applies to torsion-valued data already mapped into
$\Z/24$. It does \emph{not} place the $\Z$-valued equality/index rows of Theorem~\ref{thm:1} into $\Z/8$, and
it does not place an integer generation count into $\Z/3$. Both transfers require additional maps that are
not constructed here.

\begin{corollary}\label{cor:arenas}
\emph{Conditional torsion statement.} If an explicit map sends a sector constraint into the $\Z/8$ summand
of $\Z/24$, then that image cannot determine the separate $\Z/3$ component through the group structure alone.
This is a statement about the mapped torsion data, not an integer generation count.
\end{corollary}

Under a literal \emph{integer-index} reading, a zero or evenness constraint may instead forbid an odd count.
Thus the program-native torsion reading and the physics-default integer reading fork; neither transfers to
the other. Indeed $\Hom(\Z/3,\Z)=0$ blocks a direct class-to-count map. We retain ``selector arena'' and
``carrier arena'' only as mnemonics for the two torsion summands, not as types assigned to every obstruction
or to the generation count.

\section{Theorem 2: a finite Krein intersection invariant}\label{sec:thm2}

\paragraph{Setup (finite-dimensional).} Let $W$ be the $192$-dimensional complex generation carrier, equipped
with a nondegenerate Hermitian Krein form $K$ of signature $(96,96)$ and a chirality grading
$\Gamma\colon W\to W$, $\Gamma^2=I$, with $\pm1$-eigenspaces $W_+,W_-$ of dimension $96$ each and grading
projections $\pi_\pm=\tfrac12(I\pm\Gamma)$. The form is \emph{purely cross-chirality} (``Krein signature exactly
$(+96,-96)$,'' vectorlike; verified in \path{tests/generation-sector/ghost_parity_krein.py}): the eigenspaces
$W_+$ and $W_-$ are each $K$-isotropic (Lagrangian), $K|_{W_+\times W_+}=K|_{W_-\times W_-}=0$, and $K$ restricts
to a perfect pairing $W_+\times W_-\to\mathbb{C}$. A \emph{physical subspace} is a maximal $K$-positive-definite
subspace $P\subseteq W$ (i.e.\ $K(v,v)>0$ for all $v\in P\setminus\{0\}$); by Sylvester's law every such $P$ has
$\dim P=96$, the positive index of inertia.

\begin{definition}[Krein--chirality intersection index]\label{def:chi}
For a physical subspace $P$, set
\[
\chi_{\cap}(P):=\dim(P\cap W_+)-\dim(P\cap W_-)\in\Z.
\]
\end{definition}

\noindent This invariant counts vectors of $P$ lying \emph{entirely} in either chiral Lagrangian. It is not
the left-minus-right handedness of the $96$ physical vectors: a generic vector in a graph subspace has
components in both $W_+$ and $W_-$. A physical chirality count would require an additional
grading-invariance, graded-trace, or Fredholm/index construction.

\begin{theorem}\label{thm:2}
On the $(96,96)$ cross-chirality carrier, every physical subspace $P$ satisfies $\chi_{\cap}(P)=0$.
Consequently every
linear Krein-isometric operator $U$ (i.e.\ $K(Ux,Uy)=K(x,y)$) maps physical subspaces to physical subspaces and
preserves this nullity: $\chi_{\cap}(UP)=\chi_{\cap}(P)=0$.
\end{theorem}

\begin{proof}
Each eigenspace $W_\pm$ is $K$-isotropic, so every vector in it is $K$-null. A $K$-positive-definite
subspace contains no nonzero $K$-null vector. Hence $P\cap W_+=P\cap W_-=\{0\}$ and
$\chi_{\cap}(P)=0$. A linear $K$-isometry carries a maximal $K$-positive subspace to another such subspace,
so the same argument gives $\chi_{\cap}(UP)=0$.
\end{proof}

\begin{remark}[exact scope]
This is finite-dimensional Krein linear algebra. It does not prove a physical chiral state count, a graded
trace, spectral flow, or a Fredholm index. At canon grade, faithful stand-ins discharge the APS/end/gap/family
terms for a separately stated conditional section-setting model
(\path{canon/function-space-index-conservation-residual-closure-RESULTS.md}); transfer to the true
Rarita--Schwinger $Y^{14}$ bundle/source-action pushforward remains open.
\end{remark}

\noindent The theorem is Lean-checked in
\path{Lean/GUFormalization/LocatedNotForcedLegs.lean} using the same intersection definition. Existing numerical
scripts also report the projection-rank identity
$\dim\pi_+(P)=\dim\pi_-(P)=96$ on graph subspaces. That identity is a separate balance diagnostic and is not
counted as $96$ left states minus $96$ right states.
The numerical signature cross-check reproduces the same-chirality block nullity and zero projection-rank
difference in $(9,5)$ and $(7,7)$; a grading-aligned Euclidean $(14,0)$ control gives projection-rank
difference $96$. These are diagnostics and controls, not physical handedness indices.

\paragraph{Antilinear audit boundary.}
The finite intersection theorem does not invoke Wigner's Hilbert-space ray theorem and does not make
antilinearity a unique physical escape. A separately delimited antilinear calculation studies re-gradings whose
new chirality eigenspaces are $K$-null; for that class the same intersection index remains zero
(\path{canon/antilinear-bound-RESULTS.md};
\path{canon/antilinear-nonkrein-admissibility-RESULTS.md}). This is still Krein-side nullity, not a physical
handedness count.

\paragraph{A GU-motivated re-grading candidate is frame-trivial (computed).} The operator $C=\Jq\cdot G$ is built from
$\Jq = \id_{14}\otimes U$ (the quaternionic structure) and the chiral grading, both internal-fiber
endomorphisms. Its tangent-frame charge is exactly $0$:
$\max\|[\Jq,\ \text{any tangent-frame }\mathfrak{so}(9,5)\text{ rotation}]\| = 0.00$, including the $\Lamp$
rotation, and $|\langle \Jq,\Lamp\rangle| = 0$. This is convention-independent (any $\id_{14}\otimes U$ is
traceless on the frame factor; verified with a random $U$). Consequently this candidate carries no $p_1$, and its
reduced boundary $\eta$ is the $2$-primary gauge type $(2q^2-4q+1)/8$ (here forced to $0$ by $C^2=-I$), and the
calculation supplies no integer generation count. The denominator-eight diagnostic does not by itself define
a map from this operator to the $\Z/8$ summand of $\pis$.

\paragraph{Bounded operator audit.} The class-C and first-step-control calculations found no operator that
constructs an integer generation count. A frame-active linear overlap
$O=L_{\mathrm{SD}}\otimes X_L$ exists, so the stronger claim that frame charge and chirality support are
orthogonal is false. Its reported $+16$ and $+2$ are subspace-restriction/overlap diagnostics, not a conserved
physical chiral index. This audit is complete only for the encoded equivariant class C and the listed
controls; it does not quantify over unrestricted, non-equivariant, or function-space constructions.

For antilinear re-gradings whose new eigenspaces are $K$-null, the delimited calculation proves the same
intersection nullity as Theorem~\ref{thm:2}. Frame-nontrivial antilinear swap operators exist once symmetry is
dropped; the theorem says only that their $K$-null eigenspaces have zero intersection index. It does not
identify that index with a physical handedness count and does not close $K$-definite, unrestricted, or
true-$Y^{14}$ source-action constructions.

\paragraph{What remains open.} External chiral backgrounds are one standard way to produce an integer index:
a 2D magnetic-flux Wilson--Dirac model realizes net chiral index $=$ flux number
(\path{canon/external-topological-index-flux-RESULTS.md}). The present finite class-C audit does not prove that
external data are necessary, because it does not cover non-equivariant operators or the true
Rarita--Schwinger $Y^{14}$ bundle/source action. The supported conclusion is narrower: none of the encoded
sector-interior constructions supplies a typed integer generation count, and the missing map/operator is named
rather than inferred away.

\section{The located 3-primary carrier}\label{sec:carrier}

The $3$-primary content lives in exactly one place.\footnote{Throughout, ``order $3$'' names the $3$-primary
\emph{component} of the framed class, not the order of the class as a group element. The class is
$2\in\pis=\Z/24$, which has order $12$ as an element. Under the CRT splitting $\Z/24\cong\Z/8\oplus\Z/3$ it is
$(2\bmod 8,\,2\bmod 3)=(2,2)$: the $\Z/8$ component has order $4$, and the $\Z/3$ component $2\in\Z/3$ is the
nonzero element, of order $3$. The ``$3$-primary carrier'' is this $\Z/3$ component---the summand the $2$-primary
no-go is blind to, and the one $\Hom(\Z/3,\Z)=0$ concerns---never a claim that the element itself has order
$3$.} The self-dual $\SU(2)_+ = \Lamp$ structure, read as a
\emph{tangential} framing on $\RP^3 = L(2;1)$ (the deformation-retract spine of the metric fiber
$\mathrm{GL}(4,\R)/O(3,1)$), has framed-bordism class $\pm 2$ in $\pis = \Z/24$ and Adams $e$-invariant
\begin{align*}
\eR &= p_1/48 = (p_1/2)/24 = 1/12,\\
p_1&=4\quad\text{for the adjoint charge-$1$ bundle},\\
p_1/2&=2\quad\text{(stable framing degree; class $2$ in $\Z/24$)},\\
&\hspace{2em}\text{with nonzero $3$-primary component of order $3$.}
\end{align*}
\paragraph{Normalization conventions (fixed once, to remove generator/sign ambiguity).}
The value $\eR=1/12$ depends on exactly three conventions, which we state explicitly; given them it is rigid.
\emph{(i) Generator of $H^4(B\mathrm{Spin};\Z)$.} $H^4(B\mathrm{Spin}(n);\Z)\cong\Z$ is generated by the spin
characteristic class $\tfrac12 p_1$, so $p_1=2\cdot(\text{generator})$ (McLaughlin~\cite{mclaughlin1992};
restated by Sati--Shim~\cite{satishim2015}). The adjoint charge-$1$ ($c_2=1$) bundle has $p_1=4$ (the Dynkin
index of the $\SU(2)$ adjoint), hence $\tfrac12 p_1=2$.
\emph{(ii) Stabilization is $\times 2$.} $\pi_3(\SO(3))\cong\Z$ and $\pi_3(\SO)\cong\Z$, and the map
$\pi_3(\SO(3))\to\pi_3(\SO)$ induced by $\SO(3)\hookrightarrow\SO$ is multiplication by $2$ (standard; the
Pontryagin / framing normalization is Kirby--Melvin~\cite{kirbymelvin1999}). The charge-$1$ adjoint bundle is
the generator of $\pi_3(\SO(3))$ (via $\SU(2)=\Spin(3)\to\SO(3)$, an isomorphism on $\pi_3$); under the
isomorphism $p_1/2\colon\pi_3(\SO)\xrightarrow{\sim}\Z$ its \emph{stable framing degree} is $p_1/2=2$, not the
unstable Dynkin index $4$.
\emph{(iii) Generator of $\pis$.} $\pis\cong\Z/24=\Imm J_3$, with the real Adams $e$-invariant
$\eR\colon\pis\to\mathbb{Q}/\Z$ normalized so the generator $\nu=J(\text{generator of }\pi_3(\SO))$ has
$\eR(\nu)=1/24$ (Adams~\cite{adams1966jx4}); then $\eR$ is injective and $\eR(m\,\nu)=m/24$. Composing, the
stable framing degree $2\in\pi_3(\SO)$ maps under $J$ to the class $2\in\Z/24$, so
$\eR = (p_1/2)/24 = p_1/48 = 2/24 = 1/12$. The composite $\eR=(p_1/2)/24$ is ours, not a formula quoted verbatim
from either source. Signs are immaterial to every conclusion: replacing $\nu$ by $-\nu$ sends the class
$2\mapsto 22\equiv-2$, whose $3$-primary reduction $-2\equiv 1\pmod 3$ is again nonzero of order $3$. These are
the paper's stated conventions throughout; only the identification of GU's $\Lamp$ twist with this natural
tangential framing is reconstruction-grade.

\paragraph{It locates but does not fill (properties computed; ``locates'' is interpretation).} The carrier
$\Lamp$ has net self-dual tangent-frame charge
$33.94$ (tangential, $p_1=4$)---it is genuinely on the carrier-arena side---\emph{but it is vectorlike} (net
chiral $0$), and $\eR = 1/12$ is \emph{homotopy-fixed}: it is a fact about $\pis$, identical for a universe with
one generation or five. So the carrier marks the $3$-primary slot; it does not co-vary with the answer and so
cannot, by itself, be the thing that counts to three. The negative control confirms the slot is unique: $\RP^3$'s
own deck group is $\Z/2$, and the charge-$q$ Dirac $\eta = (2q^2-4q+1)/8$ is $2$-primary for every integer $q$,
so the $3$-primary burden sits \emph{only} in the gravitational framing channel $-p_1/24$. Throughout, ``locates
the slot'' is an interpretive gloss under the torsion-count premise (Corollary~\ref{cor:arenas}), not an
additional computed quantity: the computed facts are the tangentiality ($p_1=4$), the vectorlikeness (net chiral
$0$), and the homotopy-fixedness; and by $\Hom(\Z/3,\Z)=0$ (Section~\ref{sec:open}) the located class cannot
itself \emph{be} an integer count.

The computed picture is therefore limited but useful: the frame-trivial re-grading candidate and the
frame-charged class are distinct objects, and the latter has a nonzero $3$-primary projection. Neither object
is a typed integer generation count, and no direct class-to-count map exists.

\section{The deciding computation, and the gate}\label{sec:decider}

A direct candidate computation that could produce an integer count on GU's \emph{actual} geometry, rather than a
homotopy-fixed class identical for any answer, is the net chiral index as the anomaly-inflow coefficient of the
bulk gravitational $-p_1/24$ SPT class on the $14$-manifold, with the tangential-vs-gauge fork resolved by the
computation. We constructed and ran it ($26/26$ checks across four scripts plus an independent from-scratch
re-verification, all controls reproduced). The result:
\begin{itemize}[leftmargin=1.6em]
\item \textbf{The literal net-chiral integer is gated} on the unbuilt stabilized twisted Rarita--Schwinger / IG
source action (the ``$+8$'' Rarita--Schwinger leg). Every analytic route to it failed (ten Atiyah--Singer routes
gave $\{960,-288,-384,-192,-336,-128,128,-8,-480,60\}$, none equal to $16$); the computation explicitly refused
it rather than fabricating a fit.
\item \textbf{The only unconditionally computable generation integer in this program is one.} GU's verified Pati--Salam
$\Spin(7,7)\to\Spin(6)\times\Spin(4)$ chain gives $16$ chiral states $=$ exactly one anomaly-free generation
($\mathrm{Tr}\,Y = \mathrm{Tr}\,Q = 0$, $16/16 = 1$). The $\Spin(7,7)$ Pati--Salam chain is one reduction among
many possible symmetry-breaking chains; the integer $1$ it yields is chain-relative, not chain-independent. The
tested integer outputs are
$\{0~(\text{linear net-chiral, vectorlike}),\ 1~(\text{Pati--Salam}),\ 2~(\Ahat(K3))\}$; none is $3$.
\item \textbf{The tested re-grading candidate has a denominator-eight boundary diagnostic}
($e=3/8$, $3$-part zero) and is frame-trivial. This does not define an integer count or a canonical map into
the $\Z/8$ summand.
\end{itemize}

\paragraph{The forcing-slot test (computed).} The gate above has been probed directly. We reverse-engineered the
necessary conditions any source action must meet to force a count: a term simultaneously (a) tangential (carries
$p_1$), (b) net-chiral, and (c) non-frame-trivial. A toy stabilized twisted Rarita--Schwinger sector built four
ways (Faddeev--Popov stabilization; the twisted spin-$3/2$ Alvarez-Gaum\'e--Witten index~\cite{alvarezgaume1984grav} on $K3$; the
RS frame-index operator on the $\Cl(9,5)$ substrate; the gravitino anomaly polynomial), each angle adversarially
re-verified, reaches at most two of the three properties, never all three with a derived generation integer.
The tested integers are even or equal one: $256 = 2^8$ (a tautological projector trace, with
physical-sector index $0$), $-672 = -2^5\cdot 3\cdot 7$, $-42 \equiv 0 \pmod 3$ (the $\Z/3$ identity), the
$\mathbb{HP}^2$ unit $1$. The gravitino anomaly is not an odd-prime exception (every spin-$3/2$ coefficient is
$2$-primary up to von Staudt--Clausen~\cite{vonstaudtclausen1840} denominators), and the
twisted-by-$\mathbf{16}$ index $16(-42)+3\,\ch_2(V)$ is $\equiv 0 \pmod 3$ for every integer twist. These are
four bounded negative tests, not a classification of source actions. Their integer outputs do not acquire a
$\Z/8$ type merely because some are even. The actual stabilized action remains unbuilt.

\paragraph{The ``$2+1$'' reading is numerology across frameworks (computed).} The ``$2+1$ effective'' picture
(the $2$ from $\Ahat(K3)$, the $1$ from Pati--Salam) is not a single legitimate count. In one common framework
(the $\Ahat$-genus character-valued index on $K3$), the spin-$1/2$ leg gives $2$ but the spin-$3/2$ leg gives
$-42\equiv 0\pmod 3$, not a clean $+1$; the $+1$ appears only as the twisted-Dirac unit on a different manifold
($\mathbb{HP}^2$). The two summands of $2+1$ live in disjoint frameworks, so their sum is a coincidence, not a
derivation. ``Located, not forced'' thus holds under both the order-$3$/triality reading and the additive $2+1$
reading.

\paragraph{Carrier-mass test (computed).} In the tested finite model the carrier $\Lamp$ is vectorlike, so the
source action's possible \emph{Dirac mass} is a relevant unresolved term. The mass is allowed, not forbidden:
the Kramers structure $C^2=-I$ is pseudoreal, hence
self-conjugate, hence vectorlike, which admits a mass; and the built Seiberg--Witten action realizes a vectorlike
one. Both branches give zero, not three: massive, the $96$ generation modes pair with $96$ mirror modes and
decouple, leaving net chiral index $(n_+ - r) - (n_- - r) = 0$ analytically for every mass; massless, the $192$
stay light but the tested balance remains $0$. A light chiral count would require a separately typed
projection/index construction that breaks this balance. The finite Krein theorem does not prove that no such
physical or function-space construction exists. Consistently, exact
$\Spin(9,5)$ representation theory (computed by decomposing $S^+\otimes S^+$ over the explicit
$\mathrm{Cl}(9,5)=M(64,\mathbb{H})$ gamma matrices into the Clifford form-degree basis
$\mathrm{End}(S)=\bigoplus_k\Lambda^k$ and solving for the $\Spin(9,5)$-invariant bilinear space by
trace-orthonormal projection onto antisymmetrized Clifford words, checksum $\sum_k\mathrm{mult}=128^2=16384$;
\path{tests/chase/MOVE-4/move4_spinor_square_forms.py}, with the independent re-check
\texttt{tests/chase/MOVE-4/verify/indep\_check.py}) gives
$\dim\Hom_{\Spin(9,5)}(S^+\otimes S^+,\Lambda^0)=0$: no
invariant same-chirality Majorana scalar-mass bilinear exists in the equivariant family (the scalar bilinear
pairs only $S^+\leftrightarrow S^-$), so that mass channel is absent from the tested equivariant family. This
is a consistency remark about the spinor-square channel, not a statement about the $192$-dimensional
carrier's Dirac mass. (This Hom-vanishing is now canon as fact~(A) of the two-arena rep-theory core, \texttt{canon/two-arena-rep-theory-core-RESULTS.md}, re-run exit~0 on 2026-07-03.)

\begin{remark}[Global anomaly layer, 2026-07-03, exploration-grade]
The anomaly-inflow bridge invoked above can be probed directly on the global (Dai--Freed / spin-bordism) layer: requiring the whole anomaly-free Standard Model as chiral boundary data, does it pin a $3$-primary torsion datum? For the actual SM gauge group the tested route does not---the relevant global-anomaly group has no 3-torsion ($\Omega^{\mathrm{Spin}}_5(BG_{\mathrm{SM}})\otimes\Z_{(3)}=0$), because 3-locally $(SU(3)\times U(1))/\Z_3 = U(3)$ is torsion-free and $\Omega^{\mathrm{Spin}}_*(\mathrm{pt})$ is odd-torsion-free (Anderson--Brown--Peterson / Wall). This removes one tested mechanism; it neither defines an integer count nor proves all routes absent.
\end{remark}

\paragraph{Gated, not fabricated.} Four things remain genuinely gated: (i) the $+8$ RS leg / the literal
integer, on the unbuilt source action; (ii) the full analytic Bismut--Cheeger~\cite{bismut1989eta}
fibered-boundary reduction theorem
for the non-product $9$-dimensional $S^6$-bundle over $\RP^3$; (iii) the families pushforward
$\pi_!\colon \ch(S)/Y^{14}\to\ch(S_X)/X^4$, undefined because the fiber $\mathrm{GL}(4,\R)/O(3,1)$ is non-convex;
(iv) $3\text{-primary torsion component}\to\text{integer }3$ (Section~\ref{sec:open}). The program has previously caught
four fabricated paths to three (a disguised $\chi$, a reverse-engineered $+8$, a circular rank-$4$, a fitted
holonomy); this computation adds none, and asserts the absence ($\texttt{integer\_is\_3 = False}$).

\section{What is not claimed: the open conjecture}\label{sec:open}

We do \emph{not} claim three chiral generations. The supported verdict is \textbf{located, not forced.} The single
open bridge is
\[
3\text{-primary torsion component} \longrightarrow \text{integer }3.
\]
A nonzero class in $\Z/3$ detects information mod $3$; it is not equality to the integer $3$, and the
$e$-invariant carrying it ($\eR=1/12$) is identical for any generation count. APS, Dai--Freed~\cite{dai1994eta},
Callan--Harvey~\cite{callan1985anomalies}, and Bismut--Cheeger~\cite{bismut1989eta} relate boundary invariants
to indices, determinant lines, anomaly phases, and cobordism classes, not directly to an integer family count.
A direct class-to-count homomorphism is impossible because $\Hom(\Z/3,\Z)=0$; likewise
$\Hom(\Z/24,\Z)=0$.

An integer count, if one exists, therefore needs a separately constructed integer-valued invariant---such as a
relative, equivariant, rank, or Fredholm index---plus an explicit relation between its mod-$3$ reduction and
the located torsion component. On GU's actual $14$-manifold that would require, at minimum, a proved
fibered-boundary reduction, an explicit stabilized twisted Rarita--Schwinger operator/source action, and a
well-typed integer extraction. Those objects are not presently built.

The computations on the $\RP^3=L(2;1)$ spine and their $L(3;1)$ controls are evidence about specific
constructions, not an exhaustion theorem. Likewise, the class-C census and the finite Krein theorem do not
exclude non-equivariant operators, $K$-definite constructions, or the true-$Y^{14}$ family pushforward. The
open question is therefore narrow but genuine: can a geometry-dependent integer index be constructed on the
actual bundle, and can its mod-$3$ reduction be proved to match the exhibited torsion component? Until then,
no computation in this program yields the integer three.

\section{Relation to Geometric Unity and to prior art}\label{sec:relation}

\paragraph{Geometric Unity.} The sector is motivated by GU's matter proposal~\cite{weinstein2021gu} (the
$14$-dimensional observer space $Y^{14} = \mathrm{Met}(X^4)$, the vector-spinor field, the $4+10$ split, the
$\mathbf{16}$ of $\Spin(10)$, the self-dual base structure). The public GU material presents the fermion sector
schematically, does not build the matter action (the draft's own candidate source term, its eq.~10.10, is
author-disclaimed: ``until it is stabilized. Caveat Emptor.''), states chirality is only \emph{effective}, and
proposes a ``$2+1$'' picture rather than three true generations; the Nguyen--Polya
critique~\cite{nguyen2021response}
argues central technical details are unverifiable from the public material. Our
results are therefore a reconstruction, and where GU's own narrative points (its $2+1$ hedge---distinct from the
computed $\Ahat(K3)+\text{Pati--Salam}$ $2+1$ that Section~\ref{sec:decider} shows to be numerology across
disjoint frameworks; both, however, point away from a clean internal three---and the verified
$\Spin(7,7)\to 1$ chain), it tilts against the strong three-generation reading.

\paragraph{Prior art.} Deriving or constraining the generation number from topology or anomalies has substantial
precedent: Dobrescu--Poppitz~\cite{dobrescu2001number}; Evans, Ibe, Kehayias and
Yanagida~\cite{evans2012nonanomalous}; Kaplan--Sun~\cite{kaplan2012spacetime};
Garc\'ia-Etxebarria--Montero's $\Z_9$ anomaly forcing $N_{\mathrm{gen}}\in 3\Z$~\cite{garciaetxebarria2019daifreed};
Juven Wang's 2023 family puzzle via framed/string bordism $\Z/24$ with chiral central charge
$c_-=24$~\cite{wang2023family}; and Wan--Wang--Yau v2~\cite{wan2026fermion}, which explicitly separates
$2$- and $3$-primary anomaly data and combines the corresponding extensions by CRT. Their conditional minimal
$N=N_c=N_f=3$ conclusion also uses the generalized Standard Model without sterile $\nu_R$, minimal cyclic
extensions, the $\SU(2)$ oddness condition, color-center matching, and fermionic baryons. Deriving three
generations from algebraic structure is likewise an active lane---e.g.\ the division-algebra trialities of
Furey--Hughes~\cite{furey2024trialities}, where two generations arise as spinors and a third through a Cartan
factorization; our Clifford--Rarita--Schwinger route is mechanically distinct, but the question is not untouched.
We claim no novelty for ``topology constrains family number.'' Wang's framed/string bordism $\Z/24$ is itself $\Imm
J_3$~\cite{adams1966jx4}, and his title arithmetic $24/8 = 3$ already pulls the odd $3$-primary factor out of
$\Z/24$ as the family number, supported by a separate $w_1$--$p_1$ bordism $\Z_3$ class; Wan--Wang--Yau v2
already treats the $2/3$-primary split and CRT. We therefore claim no novelty for the factorization, primary
decomposition, CRT step, or use of $\pis=\Z/24$. To the best of the bounded primary-source search recorded in
\path{papers/candidates/located-not-forced/PRIOR-ART-DELTA.md}, no located source combines this explicit
cross-chirality Clifford--Rarita--Schwinger carrier, the scoped class-C constraint audit, and the inverse
observation that a mapped $\Z/8$ torsion image cannot determine the separate $\Z/3$ component. This is only a
search-bounded candidate delta, not proof of publication novelty, and the inverse step is elementary once the
typed premises are supplied. The finite Krein intersection lemma is likewise elementary; indefinite-metric
Standard-Model formulations with an explicit generation space are prior art~\cite{besnard2021lorentzian}.

\section{Status of claims}\label{sec:status}

{\small
\begin{longtable}{>{\raggedright\arraybackslash}p{0.34\textwidth} >{\raggedright\arraybackslash}p{0.60\textwidth}}
\toprule
\textbf{Claim} & \textbf{Grade} \\
\midrule
\endhead
Encoded class-C constraints do not supply an odd-prime congruence or select an odd integer & computed carrier census plus Lean-verified finite closure; finite torsion and $\Z$-valued rows kept in distinct codomains; no completeness outside class C \\
CRT decomposition $\pis=\Z/8\oplus\Z/3$ & standard; disjointness applies only after an explicit torsion-valued map into $\Z/24$ \\
Every maximal $K$-positive subspace has zero chiral-Lagrangian intersection index & finite-dimensional theorem, Lean-checked; not a physical handedness count or Fredholm index \\
Antilinear $K$-null re-gradings retain zero intersection index & delimited theorem; no Wigner-based uniqueness claim and no transfer to $K$-definite or unrestricted constructions \\
Tangential $\Lamp$ framing carries $\eR=1/12$ & full class has order $12$; projected $3$-primary component has order $3$; GU identification reconstruction-grade \\
The framed class is homotopy-fixed and does not define an integer count & computed/standard-result-applied; $\Hom(\Z/3,\Z)=0$ blocks a direct class-to-count homomorphism \\
Toy source-action constructions supply the missing integer count & NO in four tested constructions; bounded negative evidence, not a universal theorem \\
Carrier Dirac mass & allowed/vectorlike in the tested model; neither massive nor massless branch derives three \\
Carrier-specific inverse assembly & search-bounded candidate delta only; primary decomposition and CRT are prior art (Wang; Wan--Wang--Yau v2) \\
The literal generation integer on GU's $14$-manifold & gated on the unbuilt true-$Y^{14}$ source action/pushforward \\
The only unconditionally computable generation integer in this program & $1$ (Pati--Salam $\Spin(7,7)$, chain-relative) \\
$3$-primary torsion component $\to$ integer $3$ & open and ill-typed as a homomorphism; needs a separately constructed integer/rank/index home \\
External backgrounds are necessary & not proved; a flux-index model is a positive control, while non-equivariant and true-$Y^{14}$ routes remain open \\
GU forces exactly three chiral generations & not claimed; no computed quantity in this program equals three \\
\bottomrule
\end{longtable}
}

\section{Conclusion}\label{sec:conclusion}

The release-surviving result is narrower than the original structural reading. In the encoded equivariant
class C, finite torsion constraints are $2$-primary and integer constraints carry only parity/divisibility
content; none supplies a mod-$3$ congruence or selects the integer three. The CRT decomposition separates
torsion data only after a map into $\Z/24$ is supplied. It does not convert a $\Z$-valued index into
$\Z/8$ or a generation count into $\Z/3$.

The finite Krein theorem proves zero intersection with the chiral Lagrangians, not a physical left-minus-right
state count. The framed class $2\in\Z/24$ has order $12$ and a nonzero order-$3$ projection; it is tangential,
homotopy-fixed, and not an integer count. The true Rarita--Schwinger $Y^{14}$ source-action pushforward,
non-equivariant constructions, and a typed integer index remain open. Thus ``located, not forced'' means only
that a $3$-primary torsion component is exhibited while no tested construction supplies the missing integer
bridge. We do not claim three generations, a universal no-go, or that external backgrounds are necessary.
External review of the scoped theorems and the open typing bridge is requested.

\section*{Acknowledgments}

This work was produced by an AI-directed research process operating under the discipline of compute, then
adversarially verify, then keep only what survives. The campaign repeatedly killed its own stronger readings
under adversarial pressure (including four fabricated paths to ``three'' that were caught and rejected, and
several overclaims caught by an internal hostile-referee pass before submission); that self-correction is part of
the contribution. Geometric Unity, due to Eric Weinstein, served as the generative test case that pointed at the
structure.

\section*{Data Availability}

All load-bearing computations are reproducible from the public repository at
\url{https://github.com/disruptionjoe/gu-formalization}. A single self-contained script,
\path{papers/candidates/located-not-forced/reproduce_all.py}, recomputes the $31$ enumerated numerical and
symbolic checks in the release manifest and checks each against its stated value---exit~0 iff all match, and every check
ships a discriminating control that must fail on a scrambled input---so a referee can re-derive the paper's
numbers in one command (see \texttt{REVIEWER.md}). The representation-theoretic decompositions, the Krein
signature, the index-conservation check, the frame-charge computations, the boundary $e$-invariant controls, the
decider integers (reproduced by an independent from-scratch script, $26/26$ checks), the forcing-slot
test, and the carrier-mass capstone are checked numerically under
\texttt{tests/} (notably \texttt{tests/generation-sector/}, \texttt{tests/decider/},
\texttt{tests/boundary-eta/}, \texttt{tests/forcing-slot/}, \texttt{tests/carrier-mass/},
\texttt{tests/gu-independent/}). The homotopy and $e$-invariant facts are standard
(Adams~\cite{adams1966jx4}; Kirby--Melvin~\cite{kirbymelvin1999}). Per-result derivations are recorded in the
repository \texttt{canon/} directory, including:
\path{canon/core-theorems-symbolic-proof-RESULTS.md},
\path{canon/external-by-structure-synthesis-RESULTS.md},
\path{canon/external-topological-index-flux-RESULTS.md}, and
\texttt{canon/rs-function-space-framework-SPEC.md}.

\paragraph{Verification status.} By the project's own three-tier vocabulary, every result in this paper is at
most \emph{internally established}: computed, independently re-derived from scratch, and adversarially reviewed
\emph{within the same AI-directed process that produced it} (reproduced, not replicated). No result here is
\emph{externally established}---independently replicated, peer-reviewed, or signed off by a named specialist.
Because internal reviewers are spawned by the same process that produced the results, no internal step, however
adversarial, can cross that boundary; the request for external review is a request to cross it.

\appendix

\section{Reproducibility: the Pati--Salam one-generation verification}\label{app:patisalam}

The ``only unconditionally computable generation integer in this program is one'' of Section~\ref{sec:decider} rests on a self-contained
group-theory verification of GU's Pati--Salam chain, reproduced here for completeness. The script
\path{lab/active-research/pati_salam_chain_verification.py} (pure \texttt{numpy}) builds the $\Spin(10)$
chiral spinor $\mathbf{16}$ from its weights $(\pm\tfrac12)^5$ with an even number of minus signs, embeds
Pati--Salam by partitioning the rank-$5$ Cartan as $3\,(\text{color}+B{-}L)\oplus 2\,(T_{3L},T_{3R})$, and
computes hypercharge and electric charge from group theory alone via
\[
Y = T_{3R} + \tfrac12(B{-}L), \qquad Q = T_{3L} + Y, \qquad n = 6Y .
\]
Pushing all $16$ weights through the branching
\begin{align*}
\Spin(7,7)&\to\Spin(1,3)\times\Spin(6,4)\\
&\to\Spin(1,3)\times\Spin(6)\times\Spin(4)\\
&\cong \mathrm{SL}(2,\mathbb{C})\times\SU(4)\times\SU(2)\times\SU(2)\\
&\to \mathrm{SL}(2,\mathbb{C})\times\SU(3)\times\SU(2)\times U(1)
\end{align*}
collapses the $\mathbf{16}$ to exactly one Standard-Model generation: the multiplets $(\mathbf{3},2)_{n=+1}$,
$(\overline{\mathbf{3}},1)_{n=+2}$, $(\overline{\mathbf{3}},1)_{n=-4}$, $(\mathbf{1},2)_{n=-3}$,
$(\mathbf{1},1)_{n=+6}$, $(\mathbf{1},1)_{n=0}$, totalling $6+3+3+2+1+1=16$ states, with
$\mathrm{Tr}\,Y=\mathrm{Tr}\,Q=0$ (anomaly-free) and electric charges $\{0,\pm\tfrac13,\pm\tfrac23,\pm1\}$. The
integer label $n=6Y$ reproduces the GU draft's Section~11.3 table exactly; an embedding-ambiguity probe confirms
that the naive $B{-}L$-only hypercharge fails, so the standard $Y=T_{3R}+\tfrac12(B{-}L)$ embedding is forced. An
independent explicit-Clifford cross-check (\path{lab/active-research/verify_clifford_explicit.py}, with
$32\times32$ gamma matrices satisfying the $\SO(10)$ Clifford relations numerically) reproduces the same content
and its exact CP conjugate ($n\to-n$, $\mathbf{3}\leftrightarrow\overline{\mathbf{3}}$).

\paragraph{Scope: this is a one-generation \emph{structural} verification, not a generation count.} What is
verified is that the $\mathbf{16}$ of $\Spin(10)$ carries exactly \emph{one} anomaly-free Standard-Model
generation, with the paper's hypercharge assignments reproduced from first principles---i.e.\ internal
representation-theoretic consistency of GU's group theory. It does \emph{not} provide the generation count: it
does not establish that nature realizes this breaking pattern, does not address the effective-chirality or
``$2+1$'' claims, and---being a single generation---says nothing about \emph{how many} generations occur.
Moreover the $\Spin(7,7)$ Pati--Salam chain is one reduction among many possible symmetry-breaking chains; the
``$1$'' it yields is chain-relative, not chain-independent. It is
the integer ``$1$'' of Section~\ref{sec:decider}, and its role in this paper is precisely that: the only unconditionally
computable generation integer in this program, not a derivation of three.

\begin{thebibliography}{99}

\bibitem{adams1966jx4}
J.~F. Adams, \emph{On the groups $J(X)$. IV}, Topology \textbf{5} (1966) 21--71.
DOI:10.1016/0040-9383(66)90004-8; correction, Topology \textbf{7} (1968) 331,
DOI:10.1016/0040-9383(68)90010-4.

\bibitem{altland1997nonstandard}
A.~Altland and M.~R. Zirnbauer, \emph{Nonstandard symmetry classes in mesoscopic normal-superconducting hybrid
structures}, Phys. Rev. B \textbf{55} (1997) 1142--1161. arXiv:cond-mat/9602137.

\bibitem{alvarezgaume1984grav}
L.~Alvarez-Gaum\'e and E.~Witten, \emph{Gravitational anomalies}, Nucl. Phys. B \textbf{234} (1984) 269--330;
erratum \textbf{244} (1984) 421.

\bibitem{bismut1989eta}
J.-M. Bismut and J.~Cheeger, \emph{$\eta$-invariants and their adiabatic limits}, J. Amer. Math. Soc.
\textbf{2} (1989) 33--70.

\bibitem{besnard2021lorentzian}
F.~Besnard and C.~Brouder, \emph{Noncommutative geometry, the Lorentzian Standard Model, and its B-L
extension}, Phys. Rev. D \textbf{103} (2021) 035003. arXiv:2010.04960.
DOI:10.1103/PhysRevD.103.035003.

\bibitem{callan1985anomalies}
C.~G. Callan, Jr. and J.~A. Harvey, \emph{Anomalies and fermion zero modes on strings and domain walls}, Nucl.
Phys. B \textbf{250} (1985) 427--436.

\bibitem{dai1994eta}
X.~Dai and D.~S. Freed, \emph{$\eta$-invariants and determinant lines}, J. Math. Phys. \textbf{35} (1994)
5155--5194. arXiv:hep-th/9405012. DOI:10.1063/1.530747; erratum, J. Math. Phys. \textbf{42} (2001)
2343--2344.

\bibitem{dobrescu2001number}
B.~A. Dobrescu and E.~Poppitz, \emph{Number of fermion generations derived from anomaly cancellation}, Phys.
Rev. Lett. \textbf{87} (2001) 031801. arXiv:hep-ph/0102010.

\bibitem{evans2012nonanomalous}
J.~L. Evans, M.~Ibe, J.~Kehayias and T.~T. Yanagida, \emph{Non-anomalous discrete $R$-symmetry decrees three
generations}, Phys. Rev. Lett. \textbf{109} (2012) 181801. arXiv:1111.2481.

\bibitem{furey2024trialities}
N.~Furey and M.~J. Hughes, \emph{Three generations and a trio of trialities}, Phys. Lett. B (2025)
139473. arXiv:2409.17948. DOI:10.1016/j.physletb.2025.139473.

\bibitem{garciaetxebarria2019daifreed}
I.~Garc\'ia-Etxebarria and M.~Montero, \emph{Dai-Freed anomalies in particle physics}, JHEP \textbf{08} (2019)
003. arXiv:1808.00009.

\bibitem{kaplan2012spacetime}
D.~B. Kaplan and S.~Sun, \emph{Spacetime as a topological insulator: mechanism for the origin of the fermion
generations}, Phys. Rev. Lett. \textbf{108} (2012) 181807. arXiv:1112.0302.

\bibitem{kirbymelvin1999}
R.~Kirby and P.~Melvin, \emph{Canonical framings for 3-manifolds}, Turkish J. Math. \textbf{23} (1999) 89--115.
arXiv:math/9903056.

\bibitem{mclaughlin1992}
D.~A. McLaughlin, \emph{Orientation and string structures on loop space}, Pacific J. Math. \textbf{155} (1992)
143--156.

\bibitem{nguyen2021response}
T.~Nguyen and T.~Polya, \emph{A response to Geometric Unity}, self-published manuscript, February 23, 2021.
\url{https://files.timothynguyen.org/geometric_unity.pdf}. Not on arXiv.

\bibitem{rokhlin1952}
V.~A. Rokhlin, \emph{New results in the theory of four-dimensional manifolds}, Doklady Akad. Nauk SSSR
\textbf{84} (1952) 221--224.

\bibitem{satishim2015}
H.~Sati and H.-b. Shim, \emph{String structures associated to indefinite Lie groups}, J. Geom. Phys.
\textbf{140} (2019) 246--264. arXiv:1504.02088. DOI:10.1016/j.geomphys.2019.02.002.

\bibitem{vonstaudtclausen1840}
K.~G.~C. von Staudt, \emph{Beweis eines Lehrsatzes, die Bernoullischen Zahlen betreffend}, J. Reine Angew. Math.
\textbf{21} (1840) 372--374; T.~Clausen, Astron. Nachr. \textbf{17} (1840) 351--352.

\bibitem{wan2026fermion}
Z.~Wan, J.~Wang and S.-T. Yau, \emph{Fermion Families and Pontryagin Class: Topological Field Theory via Colour
Symmetry Extension}, arXiv:2605.26202v2 [hep-th] (2026).

\bibitem{wang2023family}
J.~Wang, \emph{Family Puzzle, Framing Topology, $c_-=24$ and $3(E_8)_1$ Conformal Field Theories:
$48/16 = 45/15 = 24/8 = 3$}, arXiv:2312.14928 [hep-th] (2023).

\bibitem{weinstein2021gu}
E.~Weinstein, \emph{Geometric Unity (author's working draft)}, self-published manuscript, April 1, 2021.
\url{https://geometricunity.org/}. Not on arXiv.

\bibitem{wigner1931gruppentheorie}
E.~P. Wigner, \emph{Group Theory and Its Application to the Quantum Mechanics of Atomic Spectra}, Academic Press,
New York, 1959 (transl. of \emph{Gruppentheorie und ihre Anwendung auf die Quantenmechanik der Atomspektren},
Vieweg, Braunschweig, 1931).

\end{thebibliography}

\end{document}
